\chapter{Methodological foundations}
\label{chap:preliminares}

\section{Purpose of the manual}
Operations research builds and analyzes quantitative models to support decisions
under constraints. This manual combines foundations with practical use: each
example turns a problem statement into QSB data and explains what the result
means. The problems are original teaching examples, with normalized inputs
available in \texttt{manual/examples/}; they do not require proprietary manuals
or files from earlier programs.

You can start with the application guide and a linear programming example, or
go directly to the chapter for your problem family. Before running a model,
read its assumptions: two models using the same words may represent different
systems. Do not confuse deterministic inventory with random demand, an acyclic
route with a general network, or a local optimum with a global one.

\section{The modeling cycle}
A rigorous study records:
\begin{enumerate}
 \item the real system, time horizon, and decision to be supported;
 \item controllable variables, parameters, and units;
 \item assumptions, uncertainty, and the scope of the data;
 \item objective, constraints, and initial conditions;
 \item solution method, backend, and execution options;
 \item feasibility checks and interpretation;
 \item sensitivity, risks, and limits of application.
\end{enumerate}
Write down the unit of every parameter. If demand is expressed per year and
lead time in days, convert them to a common basis before calculating. A number
without a unit can produce a mathematically consistent solution that is
operationally incorrect.

\section{Validity and exactness}
Validation checks structure, dimensions, fields, and family-specific rules.
It guarantees neither feasibility nor appropriate assumptions. A solver may
prove optimality for a model, return a heuristic solution, or stop at a limit.
Consult the status, algorithm, exactness, and execution notes.

In the metadata, \texttt{exact} characterizes the procedure for the supported
problem; it does not imply symbolic arithmetic or infinite precision.
\texttt{fixtureScale} identifies an educational method intended for small cases.
Numerical solution uses tolerances; compare results within a margin consistent
with their scale rather than requiring decimal strings to match exactly.
Repeatability does not prove validity either: a formulation error can be
perfectly deterministic.

\section{Computational traceability}
To reproduce an analysis, preserve the input JSON, QSB version, command, backend,
options, and result. Also record any transformations from a spreadsheet. For
external data, include the source, license, date, and unit of analysis. A
solution document does not always contain the complete model; save it separately.

The commands in this manual run from the repository root. \texttt{swift run qsb}
compiles if necessary; \texttt{.build/debug/qsb} uses the built executable. Small
examples let you check results by hand before studying larger cases. Changing
an input requires validation and solution again, rather than merely editing a
presentation table.

\begin{observacion}
The computational result expresses the consequences of the encoded data and
assumptions. The analyst must justify those assumptions, compare the solution
with the real system, and communicate its limitations to the decision maker.
\end{observacion}

\section{Comparison and sensitivity}
Distinguish verification (the program computes what was specified) from model
validation (what was specified adequately represents the system). Check balances
and constraints using the returned values. Then vary demand, costs, or
probabilities while keeping their distributions and units consistent. A change
in recommendation after a small perturbation may reveal a fragile decision; a
tie may require secondary criteria outside the model. Do not extrapolate a
conclusion beyond the range analyzed.
